\section{Later-Confirmed Self-Improvement and Depth-First AGA Evidence}
\label{sec:aion-aga-depth-retention}

This development increment advanced the Autonomous General Apprentice (AGA)
evidence registry from three to five earned gates while preserving a fail-closed
claim boundary.  The newly earned gates were project volume and later-confirmed
cognitive-code improvement.  The resulting evidence does not authorize an AGA
claim: practical depth, expert competence, a third naturally occurring repair,
week-scale retention, and human-grounded social and creative competence remain
open.

\subsection{Later-confirmed cognitive-code self-improvement}

AION detected a naturally occurring weakness in its executive routing policy.
Five distinct objective families had collapsed into a single method,
\texttt{research\_investigation}.  Rather than modifying the live controller
directly, AION constructed an abstract-syntax-tree-restricted Python challenger
inside a private workspace.  The incumbent champion was frozen and hashed
before experimentation.

The incumbent routed five of fourteen accumulated examples correctly, whereas
the private challenger routed all fourteen.  The challenger also retained all
five protected legacy routing behaviours and rejected six of six malicious code
candidates.  Tournament success alone did not authorize promotion.  The
challenger was first installed as a proposal-only component and was subsequently
evaluated through later operational consequences across four distinct executive
methods.  Seven later-confirmed projects exceeded the minimum requirement of
three, the original champion remained byte-identical, and no unsafe write was
performed.  HexCore therefore promoted
\texttt{procedure\_later\_confirmed\_cognitive\_routing\_code\_v1}.

The resulting improvement chain was
\[
\text{natural routing failure}
\rightarrow
\text{private code challenger}
\rightarrow
\text{protected tournament}
\rightarrow
\text{proposal-only deployment}
\rightarrow
\text{later operational consequences}
\rightarrow
\text{governed promotion}.
\]

\subsection{Useful-project volume and authority diversity}

The North-Star useful-work controller exceeded its project-volume requirement.
At the initial promotion checkpoint it had completed 21 later-confirmed projects
against a requirement of 20, spanning five of five authority families with
100\% weakest-family success, zero owner interventions, and zero unsafe writes.
Continued autonomous operation subsequently increased the portfolio to 28
later-confirmed projects.  These projects include research investigations,
software tools, data decisions, machine-checked mathematics, and exact-provenance
document work.  Project generation and delayed scoring continue while other
curriculum work is parked for elapsed retention.

\subsection{Correction of the progressive competency controller}

An audit identified a mismatch between the progressive competency assessor and
curriculum generator.  The assessor correctly required unfamiliar projects,
independent outcomes, sufficient trial volume, low scaffolding, debugging,
transfer, and delayed retention.  The generator did not explicitly schedule all
of these requirements and could move prematurely toward retention after
accumulating several familiar projects.

The generator was corrected to schedule each requirement independently:
\begin{enumerate}
    \item repeated knowledge and practical evidence for every declared subskill;
    \item the required total project count;
    \item a separate unfamiliar-project threshold;
    \item debugging recovery and source-disjoint transfer;
    \item delayed closed-book retention;
    \item independently owned outcomes, sufficient trials, and declining
          scaffolding; and
    \item fresh consistency trials only after all preceding gates are satisfied.
\end{enumerate}

Unexecuted retention appointments created under the previous policy are retained
for audit as \texttt{superseded\_missing\_prerequisite}.  They award no evidence,
and no historical evidence is deleted or reclassified.

\subsection{Six-family depth acceleration}

A depth-first campaign now targets one competency from each of six qualitatively
different families, as shown in Table~\ref{tab:aion-depth-targets}.

\begin{table}[htbp]
\centering
\caption{Cross-domain Advanced targets.}
\label{tab:aion-depth-targets}
\begin{tabular}{ll}
\hline
\textbf{Family} & \textbf{Target competency} \\
\hline
Computing & Algorithms and Data Structures \\
Engineering & Software Engineering \\
Formal reasoning & Mathematics \\
Human language & English \\
Programming & Python \\
Science & Scientific Method \\
\hline
\end{tabular}
\end{table}

Algorithms and Data Structures and Software Engineering form the first Expert
lane after Advanced competence is independently earned.  The depth controller
changes scheduling priority only; it cannot alter competency labels or award
evidence.

To avoid returning to synthetic exercises, AION revalidated twelve completed
North-Star projects and bound them into the competency ledger: three each for
Mathematics, English, Python, and Scientific Method.  Every imported record
required a pre-action commitment, a later outcome hash, a passing independent
evaluation, at least sixty seconds of authority delay, zero owner intervention,
and an exact SHA-256 match for the retained artifact.  Re-execution imports no
duplicates.  Crucially, this bridge awards no delayed-retention credit.

The six target competencies have therefore met their immediate practical-breadth
requirements and remain Intermediate until fresh closed-book retention tests are
executed.  The scheduled assessments begin on 2 August 2026: Algorithms and Data
Structures at 17:39 CEST, Software Engineering at 18:23 CEST, and English,
Mathematics, Python, and Scientific Method at approximately 22:40 CEST.  A pass
may advance a subject to Advanced; a failure retains the Intermediate level and
triggers remediation.  Advanced status unlocks, but does not automatically
grant, the deeper Expert curriculum.

\subsection{Arena v15 CAU promotion}

Arena v15 completed fifteen committed cycles over 24.04 hours using four
independently changing public authorities.  An internal CAU reviewer reread the
state and append-only ledger from disk, recomputed every outcome digest, and
confirmed exact correspondence between all fifteen pre-action commitments and
outcome records.  Execution and transactional checks remained safe, and process
reconstruction succeeded.  The procedure
\texttt{procedure\_long\_duration\_real\_outcome\_campaign\_v15\_6fd206e71b24}
was therefore internally promoted.

This 24-hour result is intentionally separate from the seven-day AGA retention
gate.  A persistent campaign service has replaced the earlier bounded
supervisor.  It records a new independent observation approximately every two
hours and remains active until explicitly stopped.  The retention counter now
uses the observed interval between the first and latest committed outcomes;
passive calendar time while the service is stopped cannot produce evidence.

\subsection{Precommitted week-scale retention challenge}

The seven-day gate does not represent a seven-day lesson.  It tests whether the
integrated learning organism preserves commitments, evidence, operating state,
restart continuity, and protected capabilities across genuine elapsed
operation.  To prevent elapsed time alone from satisfying the gate, AION has
precommitted a protected week-end challenge with commitment prefix
\texttt{8544d2bc91ecdd57}.

The final task seed is defined by
\[
s_{7d}=\operatorname{SHA256}
\left(c_{0}\,\|\,o_{\geq 7d}\right),
\]
where $c_{0}$ is the frozen challenge commitment and $o_{\geq 7d}$ is the first
Arena outcome observed at or after the seven-day boundary.  Since that outcome
does not yet exist, neither AION nor the development process can know the final
task instances in advance.

The challenge will test a fresh mathematical identity, canonical-hash tamper
rejection, a fresh transactional constraint, counterexample rejection,
preservation of the protected evidence sets for all six depth targets, and
restart progression.  Both seven observed days and a passing challenge are now
mandatory.  The earliest temporal boundary is approximately 21:42 CEST on
7 August 2026, with recognition occurring on the first qualifying observation
after that time.

\subsection{Explicit AGA scopes and current position}

The registry now reports two non-interchangeable scopes.  \emph{Full AGA}
contains all eleven gates, including human-grounded social and creative
competence.  \emph{Technical AGA} contains ten gates and explicitly excludes
that capability.  Technical scope therefore cannot be presented as full AGA.
Neither scope is currently authorized.

\begin{table}[htbp]
\centering
\caption{AGA evidence status at this checkpoint.}
\label{tab:aion-aga-status}
\begin{tabular}{lll}
\hline
\textbf{Gate} & \textbf{Observed} & \textbf{Status} \\
\hline
Later-confirmed useful projects & $28/20$ & Passed \\
Outcome-authority families & $5/5$ & Passed \\
Measured method transfers & $3/3$ & Passed \\
Later-confirmed cognitive improvement & $1/1$ & Passed \\
Zero-owner-intervention streak & Requirement met & Passed \\
Natural self-repairs & $2/3$ & Open \\
Advanced competencies & $0/6$ & Delayed tests scheduled \\
Expert competencies & $0/2$ & Open \\
Replay-free retention & Approximately $1.05/7$ days & In progress \\
Human social/creative authority & $0/1$ & Blocked or excluded by scope \\
\hline
\end{tabular}
\end{table}

The earned total remains five of eleven gates, and AGA authorization remains
\textbf{false}.  The unchanged gate count is intentional: scheduling, imported
project breadth, internal promotion, and elapsed calendar age cannot substitute
for delayed practical outcomes.  Meanwhile, the curriculum continues on other
subjects---currently JavaScript and TypeScript---alongside North-Star projects,
natural-failure monitoring, private repair attribution, competency
consolidation, and persistent external observation.

Focused verification passed 28 of 28 tests for the depth, evidence-bridge, CAU,
registry, and dashboard stack, followed by four of four focused tests for the
strengthened week-retention integration.  No Git commit was created during this
increment.

